\documentclass{article}
\usepackage{amsmath,amssymb}
\usepackage{xurl}
\usepackage[hidelinks]{hyperref}
\setlength{\emergencystretch}{2em}
% Shared commands for notes in docs/paper-gaps/.

\DeclareUnicodeCharacter{2080}{\ifmmode{}_{0}\else\textsubscript{0}\fi}
\DeclareUnicodeCharacter{2081}{\ifmmode{}_{1}\else\textsubscript{1}\fi}
\DeclareUnicodeCharacter{2082}{\ifmmode{}_{2}\else\textsubscript{2}\fi}
\DeclareUnicodeCharacter{2099}{\ifmmode{}_{n}\else\textsubscript{n}\fi}

\newcommand{\Lean}{\textsc{Lean}}
\ExplSyntaxOn
\NewDocumentCommand{\leanid}{m}{
  \ifmmode
    \textsf{\detokenize{#1}}
  \else
    \group_begin:
    \urlstyle{sf}
    \tl_set:Nn \l_tmpa_tl {#1}
    \tl_replace_all:Nnn \l_tmpa_tl {\_} {_}
    \exp_args:NV \path \l_tmpa_tl
    \group_end:
  \fi
}
\ExplSyntaxOff
\providecommand{\tr}{\operatorname{tr}}
\newcommand{\tnleanIssueBase}{https://github.com/LionSR/TNLean/issues/}
\newcommand{\tnleanPrBase}{https://github.com/LionSR/TNLean/pull/}
\newcommand{\tnleanDocBase}{https://sirui-lu.com/TNLean/docs/}
\newcommand{\ghissue}[1]{\href{\tnleanIssueBase#1}{GitHub issue \##1}}
\newcommand{\ghpr}[1]{\href{\tnleanPrBase#1}{PR \##1}}
\newcommand{\doclink}[1]{\href{\tnleanDocBase#1.html}{\path{#1}}}

% Dirac notation and the identity operator, as in blueprint/src/macros/common.tex.
% \providecommand keeps note-local definitions authoritative.
\usepackage{dsfont}
\providecommand{\ket}[1]{|#1\rangle}
\providecommand{\bra}[1]{\langle #1|}
\providecommand{\braket}[2]{\langle #1 | #2 \rangle}
\providecommand{\ketbra}[2]{|#1\rangle\!\langle #2|}
\providecommand{\Id}{\mathds{1}}
\providecommand{\one}{\mathds{1}}
\providecommand{\C}{\mathbb{C}}
\providecommand{\MN}[1]{M_{#1}(\C)}
\providecommand{\MD}{M_D(\C)}

% Verdict marker: \gapnote{<kind>}{<status>}. Kind is the policy
% classification (clarification, local-correction, scope-restriction,
% unfaithful, false-source, open-gap); status is open, wip (elimination
% underway), resolved, or historical. Machine-read by the site index;
% prints a verdict line.
\newcommand{\gapnote}[2]{\par\noindent\textbf{Verdict:} #1 (#2).\par}

\title{Endpoint Adjoint Closure and the Mixed Operator}
\author{}
\date{2026-11-16}
\begin{document}
\maketitle
\gapnote{scope-restriction}{open}

\section{The proposed implication}
The mixed sectors of the interpolation in
\cite[Section~5]{GarreRubio2023MPOSym} contract synthesis at one
endpoint with analysis at the other. Adjoint closure of the endpoint
boundary ranges does not control those contractions. In particular,
unitarity of the physical endpoint operators does not imply adjoint
closure of the boundary range of this particular mixed operator.
This concerns the proposed elimination of the extra hypothesis in
\texttt{glm23\_mixed\_path\_parent\_commutation.tex}; it does not
exclude a different choice of virtual coordinates or a different path.
\footnote{Source traceability: \path{Papers/2203.12563/REsubmission.tex},
    lines 1600--1631, the mixed sectors preceding \texttt{Agammasym}.
    The proposed implication is issue \#10763, not a separately stated
    theorem of the paper.}

\section{Unitary involutive endpoint data}
Let the first endpoint have physical dimension four and bond dimension
two. With matrix units indexed from zero, set
\begin{align}
    S & =\begin{pmatrix}2&1\\1&1\end{pmatrix},            &
    Z & =\begin{pmatrix}1&0\\0&-1\end{pmatrix},           &
    Q & =SZS^{-1}=\begin{pmatrix}3&-4\\2&-3\end{pmatrix}.
    \label{eq:mea_virtual}
\end{align}
Define $A_0^{rs}=SE_{rs}S^{-1}$ and
$u_0=\operatorname{diag}(1,-1,-1,1)$. Then $Q^2=I$, $u_0$ is
unitary and Hermitian, and
\begin{align}
    \sum_{kl}(u_0)_{rs,kl}A_0^{kl}=QA_0^{rs}Q.
    \label{eq:mea_action}
\end{align}
Thus $V_0=W_0=Q$ give an exact action decomposition with multiplicity
one and $V_0W_0=I$. The letters are linearly independent. The second
endpoint has physical and bond dimensions one, $A_1=1$, $u_1=1$,
and $V_1=W_1=1$. Both endpoint MPO bond dimensions are one.
Their arbitrary-boundary operators are scalar multiples of
$u_p^{\otimes N}$, so both ranges are adjoint-closed at every length.

These data extend to the two-element group: the identity label has
identity physical and virtual actions. The generator squares to the
identity both physically and virtually. All fusion and action
multiplicities are one and all chosen $L$-symbols are one at both
endpoints. Consequently matching the algebraic group data does not
remove the obstruction below.

\section{The two-site obstruction}
Use mixed physical labels $(r,s)$ with $r,s\in\{0,1,2\}$ and regard
$2$ as the second endpoint. The first cross sector has the physical
matrix $Q^T$, since its coefficient at output $(r,2)$ and input
$(k,2)$ is $Q_{kr}$. Let $\widetilde T$ be the mixed MPO and put
$B_X=O^2_{\widetilde T,X}$. For the two physical words
$a=((0,0),(0,2))$ and $b=((0,0),(1,2))$, direct contraction gives
\begin{align}
    (B_Y)_{a,a} & =3Y_{10},  &
    (B_Y)_{a,b} & =2Y_{10},  &
    (B_Y)_{b,a} & =-4Y_{10}.
    \label{eq:mea_entries}
\end{align}
Choose $X=E_{10}$. If $B_Y=B_X^\dagger$, the diagonal entry in
(\ref{eq:mea_entries}) forces $Y_{10}=1$, while the $(a,b)$ entry
forces $2Y_{10}=-4$, a contradiction. Thus the mixed two-site
boundary range is not adjoint-closed. Moreover, the letters with equal
output and input labels $0,2,6,8$ are $E_{00},3E_{01},3E_{10},E_{11}$.
They span the full virtual matrix algebra, so the mixed operator is
one-site injective after its physical indices are combined.
\footnote{The square-alphabet assertion is expressed by
    \leanid{MPSTensor.MPOSymmetry.MixedEndpointAdjointObstruction.not_mixed_adjoint_closed}.
    The endpoint action, injectivity, physical unitarity and endpoint
    boundary closure are separate supporting declarations in
    \doclink{TNLean/MPS/Symmetry/MPOSymmetry/MixedPath/MPO/MixedEndpointAdjointObstruction}.}

\section{Polar compression does not remove this example}
The coefficient matrix of $A_0$, with rows $(r,s)$ and columns $(i,j)$,
is positive definite:
\begin{align}
    P=S\otimes S^{-1}
      & =\begin{pmatrix}
             2 & -2 & 1 & -1 \\-2&4&-1&2\\1&-1&1&-1\\-1&2&-1&2
         \end{pmatrix}=R^TR,\notag \\
    R & =\begin{pmatrix}
             1 & -1 & 0 & 0 \\0&1&0&0\\1&-1&1&-1\\0&1&0&1
         \end{pmatrix},      & \det R & =1.
    \label{eq:mea_polar}
\end{align}
Hence its polar isometry is the identity and its positive polar tensor
is $A_0$ itself. The scalar endpoint has the same property. The
arbitrary-alphabet construction therefore first forms exactly
$\widetilde T$, then includes its nine-dimensional physical space
isometrically into the fourteen-dimensional common physical space.
Compression back to that space takes each enlarged boundary operator
to $B_Y$ and commutes with adjunction. Adjoint closure of the enlarged
range would imply the closure refuted in (\ref{eq:mea_entries}).

The Gram identity and the square-alphabet obstruction have explicit
formal statements. The identification of this example's polar
isometry and the final compression argument in this section are
written mathematics; they are not claimed as a checked theorem about
the arbitrary-alphabet counterexample.

\section{A sufficient local adjoint criterion}
For each endpoint $p$, action multiplicity $\mu$ and MPO bond index
$\alpha$, regard
\begin{align}
    (V_{p,\mu,\alpha})_{rs} & =(V_{p,\mu})_{r,(\alpha,s)}, &
    (W_{p,\mu,\alpha})_{rs} & =(W_{p,\mu})_{(\alpha,r),s}
    \label{eq:mea_slices}
\end{align}
as square matrices on the MPS bond. Suppose primed endpoint data
satisfy physical adjoint equality $T'_p=T_p^\dagger$ without
reversing virtual order, and $V'_{p,\mu,\alpha}=V_{p,\mu,\alpha}^\dagger$,
$W'_{p,\mu,\alpha}=W_{p,\mu,\alpha}^\dagger$. Then
\begin{align}
    \widetilde T(T'_0,T'_1,V'_0,V'_1,W'_0,W'_1)
    =\widetilde T(T_0,T_1,V_0,V_1,W_0,W_1)^\dagger.
    \label{eq:mea_dual}
\end{align}
Indeed, diagonal sectors use the endpoint identities. In a cross
sector the conjugated coefficient is
$\sum_\mu\overline{(W_0)_{(\alpha,r),k}}
    \overline{(V_1)_{l,(\beta,s)}}$, which is precisely the primed
coefficient by (\ref{eq:mea_slices}).
\footnote{\leanid{MPSTensor.MPOSymmetry.physicalAdjointTensor_mixedEndpointMPO}.}

A rectangular physical change of coordinates preserves this identity,
since $(FTF^\dagger)^\dagger=FT^\dagger F^\dagger$. In the self-adjoint
case, polar compression and common-space inclusion therefore give
\begin{align}
    \bigl(O^N_{\widetilde T,X}\bigr)^\dagger
    =O^N_{\widetilde T,\overline X}
    \label{eq:mea_boundary}
\end{align}
at every length, including zero. For injective endpoint tensors with
exact action decompositions, (\ref{eq:mea_boundary}) discharges the
adjoint-closure hypothesis of the existing commutation theorem. The
mixed interaction commutes with every two-site boundary operator for
all $\gamma\in[0,1]$, and the periodic Hamiltonian commutes with the
periodic common MPO at every length at least two.
\footnote{\leanid{MPSTensor.MPOSymmetry.conjTranspose_mpoWithBoundary_arbitraryPhysicalMixedMPO},
    \leanid{MPSTensor.MPOSymmetry.arbitraryPhysicalMixedInteraction_commute_of_endpoint_adjoint}
    and
    \leanid{MPSTensor.MPOSymmetry.hamiltonian_arbitraryPhysicalMixedInteraction_commute_of_endpoint_adjoint}.}

This sufficient criterion adds conditions on the action maps, not
merely on the endpoint boundary ranges. For on-site involutive unitary
symmetries, the existing Hermitian-gauge theorem already derives closure
when each action gauge is proportional to a Hermitian matrix. It derives
that condition for injective canonical endpoint tensors whose transfer
map fixes the identity and whose adjoint transfer map has a positive
definite fixed point.\footnote{\leanid{MPSTensor.MPOSymmetry.exists_isHermitian_smul_bondOneActionGauge}
    and
    \leanid{MPSTensor.MPOSymmetry.exists_conjTranspose_mpoWithBoundary_arbitraryPhysicalMixedMPO_eq},
    in \doclink{TNLean/MPS/Symmetry/MPOSymmetry/MixedPath/MPO/ArbitraryPhysicalMixedHermitianGauge}.}
The present criterion instead permits arbitrary endpoint MPO bond
dimensions and multiplicities, with exact slice-adjoint identities,
and also describes dual pairs. Deriving suitable identities after a
general allowed virtual normalization, and treating non-self-dual labels
by a common adjoint-closed family, remain separate obligations.
Even for a unitary group representation a single label's
boundary range need not be adjoint-closed: the adjoint may belong to the
inverse label. The group family and each individual operator must be
distinguished.

\bibliographystyle{alpha}
\bibliography{references}
\end{document}
